% Content of the QECLean blueprint, shared by the web and pdf versions.
%
% The nodes themselves are NOT written here. They are extracted from Lean by
% LeanArchitect out of the annotations in QECBlueprint.lean, and pulled in by
% the \input below; this file supplies the narrative and decides the order in
% which nodes appear. To change a statement or a proof sketch, edit the
% annotation in QECBlueprint.lean and re-run `lake build QECBlueprint:blueprint`
% -- do not edit the generated files under .lake/build/blueprint.

\input{../../.lake/build/blueprint/library/QECBlueprint}

\chapter{Introduction}
\label{chap:intro}

This is the blueprint for \href{https://github.com/Stavan-Jain/QECLean}{QECLean},
an ongoing formalization effort for quantum error correction in Lean~4. It consists of an informal mathematical exposition along with a dependency graph that shows how the formalized definitions and theorems in QECLean relate to each other. 

\section{What QECLean contains}

The final goal of QECLean is to provide a comprehensive formalization of Quantum Error Correction theory. As of August 2026, we've built abstractions to formalize the proofs of the canonical $[[n, k, d]]$ parameters of stabilizer codes. 

Our formalization begins with the definition of the n-qubit Pauli group $\mathcal{P}_n$ (\cref{chap:pauli}) and builds theory leading up to the definitions of a StabilizerCode and code distance (\cref{chap:codes-abstract}). We then instantiate our abstract StabilizerCode definition with a variety of well-known quantum error correcting codes (\cref{chap:toric}, \cref{chap:bb}, \cref{chap:codes}). Along the way, we develop frameworks to prove properties of these codes with the goal of making the proofs as general and modular as possible. In \cref{chap:symplectic}, we present the binary symplectic representation to provide efficient procedures to prove commutation and independence properties of stabilizer generators. We also provide extensive support for CSS codes via the homological framework (\cref{chap:homological}). 

Our most notable formalizations are the following:

\begin{enumerate}
  \item \textbf{The toric family.} For every $L \ge 2$, we verify that the $L \times L$ toric
    code is a $[[2L^2, 2, L]]$ stabilizer code. Notably, we parametrically prove that the distance is exactly $L$.
    (\cref{thm:toric-distance}).
  \item \textbf{The gross code.} We verify IBM's gross code [Bravyi et al.] as a $[[144,12,12]]$ bivariate-bicycle code. Notably, we prove that the distance is exactly 12 (\cref{thm:gross-distance}).
\end{enumerate}

Alongside these sit the rotated surface code family, the repetition codes, the
iceberg family, and the small named codes: Steane's $[[7,1,3]]$, Shor's
$[[9,1,3]]$, the perfect $[[5,1,3]]$ code, and the $[[4,2,2]]$ detection code
(\cref{chap:codes}).

Everything on the main branch depends only on the standard axioms
\texttt{propext}, \texttt{Classical.choice}, and \texttt{Quot.sound}.

\section{Reading the graph}

An edge $A \to B$ in the dependency graph means that $B$ is the nearest
blueprint ancestor of $A$: LeanArchitect walks the type and the proof term of
each annotated declaration and stops as soon as it reaches another annotated
declaration. The edges are therefore read off the proof terms rather than
maintained by hand.

\chapter{The Pauli group}
\label{chap:pauli}

In the textbook Quantum Computation and Quantum Information, Nielsen and Chuang write ``The key to the power of the stabilizer formalism lies in the clever use of group theory". The group of interest is the Pauli group $\mathcal{P}_n$ on $n$ qubits. It is defined as the group $n$-qubit Pauli operators with phases $\{\pm 1, \pm i\}$. Our formalization splits an element of this group into a phase exponent in $\Z_4$ and a phase-free operator $\mathrm{Fin}\,n \to \{I,X,Y,Z\}$. The split makes multiplication easy: the operator part multiplies
qubitwise, and the phases that fall out of products such as $XZ = -iY$
accumulate separately. 

\inputleannode{def:pauli-operator}
\inputleannode{def:pauli-group-element}
\inputleannode{def:nqubit-pauli-operator}
\inputleannode{def:nqubit-pauli-group}

The syntactic representation of $\mathcal{P}_n$ as an inductive type is tied to actual operators on a Hilbert space by
a matrix interpretation:

\inputleannode{def:pauli-to-matrix}
\inputleannode{def:pauli-mulop}
\inputleannode{thm:pauli-mul-assoc}

\section{Weight}

The following definitions build up to and are motivated by the definition of code distance in \cref{chap:codes-abstract}:

\inputleannode{def:pauli-support}
\inputleannode{def:pauli-weight}

\section{Commutation}

To prove that a group of Paulis form a stabilizer group, we need to show that the group is abelian. A convenient fact that helps us prove that two Paulis commute is that commutation in
$\mathcal{P}_n$ is a parity condition. In particular, two Paulis pick up a sign $-1$ for each
qubit where they disagree nontrivially (i.e. the operators on the qubit are both not $I$). If they disagree nontrivially on k qubits, the global sign is $(-1)^k$. Thus, two Paulis commute if they disagree nontrivially on an even number of qubits (k is even) and anticommute if they disagree nontrivially on an odd number of qubits (k is odd).

\inputleannode{def:anticommute}
\inputleannode{def:anticommutes-at}
\inputleannode{thm:commutes-iff-even}
\inputleannode{thm:commute-or-anticommute}

\chapter{The binary symplectic representation}
\label{chap:symplectic}

This chapter introduces the binary symplectic framework that allows us to prove properties about Paulis by doing linear algebra over $\F_2$. In particular, we view Pauli operators as ``symplectic" vectors in $\F_2^{2n}$. Commutation of two Paulis is interpreted as orthogonality of their corresponding symplectic vectors and the independence of a set of paulis is interpreted as linear independence of the matrix of symplectic vectors. Thus, commutation becomes an orthogonality check and independence of generators becomes a
rank computation.

\inputleannode{def:to-symplectic}
\inputleannode{thm:to-symplectic-injective}
\inputleannode{def:symplectic-inner}
\inputleannode{thm:commutes-iff-symplectic}
\inputleannode{def:check-matrix}
\inputleannode{def:rows-linear-independent}

\chapter{Stabilizer groups and the codespace}
\label{chap:stabilizer}

This chapter defines the key abstractions of the stabilizer formalism and sets up the group theoretic properties that are needed to define logical operators in \cref{chap:codes-abstract}. 

\inputleannode{def:is-stabilized-by}
\inputleannode{def:stabilizer-group}
\inputleannode{thm:neg-identity-not-mem}
\inputleannode{def:codespace}
\inputleannode{def:stabilizer-sum}

The two conditions in \cref{def:stabilizer-group} ( abelian, and not
containing $-I$) are exactly what is needed for the codespace to be nonzero.

\inputleannode{thm:codespace-nonempty}

\section{The centralizer}

The analysis of error-correcting properties of a stabilizer code begins with an analysis of how operators behave on the code's codespace. The behavior of an operator on the codespace is directly connected to whether it commutes or anticommutes with the stabilizers. This section defines the subgroup that identifies the operators that commute with the stabilizers. Later, in \cref{thm:pauli-logical-iff-centralizer}, we show that these are exactly the Pauli operators that preserve the codespace.

\inputleannode{def:centralizer}
\inputleannode{thm:normalizer-eq-centralizer}
\inputleannode{thm:stabilizer-le-centralizer}

\chapter{Logical operators, codes, and distance}
\label{chap:codes-abstract}

The encoded information is acted on by the quotient
$\mathcal{C}(\mathcal{S})/\mathcal{S}$: operators that preserve the codespace,
modulo those that act trivially on it.

What makes an operator logical is semantic --- it must map the codespace into
itself --- and that is how the definition is set up here, first for an arbitrary
unitary and then for a Pauli. That this coincides with the syntactic condition of
commuting with every stabilizer is a theorem, not a definition. It is the reason
every logicality check later in the development can be a finite commutation test.

\inputleannode{def:logical-gate}
\inputleannode{thm:logical-gate-iff}
\inputleannode{def:pauli-logical}
\inputleannode{thm:pauli-logical-iff-centralizer}
\inputleannode{thm:logical-iff-generators}
\inputleannode{def:nontrivial-logical}
\inputleannode{thm:nontrivial-logical-iff}

\Cref{def:nontrivial-logical} has three conditions, not two, and the third is
easy to forget: no stabilizer may share the operator part of $g$. Without it a
$g$ differing from a stabilizer only by a phase would count as acting
nontrivially, and the CSS bridge argument of \cref{thm:not-both-boundary}
would fail.

\inputleannode{thm:nontrivial-logical-transfer}
\inputleannode{def:logical-qubit-ops}

\section{Packaging a code}

\inputleannode{def:generators-independent}
\inputleannode{thm:independent-of-check-matrix}
\inputleannode{def:stabilizer-code}

Instantiating \cref{def:stabilizer-code} for a concrete family \emph{is} the
theorem that the family encodes $k$ qubits into $n$: the type carries $n$ and
$k$, and the obligations --- generator count, independence, pairwise
commutation, exclusion of $-I$, and the centralizer and anticommutation
conditions on the logical operators --- must all be discharged to build an
inhabitant.

\section{Distance}

\inputleannode{def:has-code-distance}
\inputleannode{thm:distance-min-weight}
\inputleannode{def:code-with-distance}

Both halves of \cref{def:has-code-distance} matter. The lower bound is the
error-correction guarantee; the witness of exact weight is what makes the
distance a value rather than a bound, and it is usually the easier half to
forget.

\chapter{CSS structure}
\label{chap:css}

A CSS code is one whose stabilizer generators are each purely $X$-type or
purely $Z$-type. The separation is what lets a homological argument apply: the
$X$ and $Z$ sides become the two ends of a chain complex.

\inputleannode{def:x-type}
\inputleannode{def:z-type}
\inputleannode{thm:css-distance-two}

\chapter{The homological framework}
\label{chap:homological}

This chapter is the abstract heart of the development. A length-three chain
complex over $\F_2$ induces a CSS code: $1$-chains are the physical qubits,
$\partial_1 \partial_2 = 0$ is the commutation of the $X$- and $Z$-checks, and
--- the slogan --- \emph{logical operators are homology classes}. Both the
toric family and the gross code are instances, and both inherit their distance
argument from \cref{thm:chain-weight-transfer}.

\inputleannode{def:homological-code}
\inputleannode{thm:boundary-comp-zero}
\inputleannode{def:cycles}
\inputleannode{def:boundaries}
\inputleannode{thm:boundaries-le-cycles}
\inputleannode{def:H1}
\inputleannode{thm:finrank-H1}

\section{From chains to Paulis}

\inputleannode{def:chain-x-operator}
\inputleannode{def:chain-z-operator}
\inputleannode{thm:chain-x-add}
\inputleannode{thm:x-logical-iff-cycle}

\section{Distance from homology}

\inputleannode{thm:not-both-boundary}
\inputleannode{thm:chain-weight-transfer}

\Cref{thm:chain-weight-transfer} is the load-bearing step: it is what allows a
statement about minimum weights of homology representatives --- a question in
$\F_2$ linear algebra --- to settle a statement about the weights of Pauli
operators. Every concrete distance theorem below is an application of it.

\chapter{The toric code}
\label{chap:toric}

The $L \times L$ square lattice on the torus, read as a chain complex. The
$2L^2$ edges are the qubits; vertices carry the $Z$-checks and faces the
$X$-checks.

\inputleannode{def:toric-chain-complex}
\inputleannode{def:toric-H1}
\inputleannode{thm:toric-H1-dim}

\section{Wrapping invariants}

The homology of the torus is detected by two $\F_2$-valued winding numbers,
counting how often a chain wraps horizontally and vertically. They are defined
on all chains, are unchanged by adding a boundary, and together separate
homology classes.

\inputleannode{def:wrapping-h}
\inputleannode{def:wrapping-v}
\inputleannode{thm:wrapping-boundary-zero}
\inputleannode{def:toric-phi}
\inputleannode{thm:toric-phi-equiv}

\section{The code and its distance}

\inputleannode{def:toric-code}

The natural generating set of the toric code has $2L^2$ elements, two more than
the $n - k = 2L^2 - 2$ that \cref{def:stabilizer-code} demands: the vertex
checks multiply to the identity, and so do the face checks. The packaging
therefore uses a trimmed generator list, and the homological identities are
what prove the trimmed list generates the same subgroup. Distance proofs stated
against the untrimmed stabilizer group are carried across by
\cref{thm:nontrivial-logical-transfer}.

\inputleannode{thm:toric-dX}
\inputleannode{thm:toric-dZ}
\inputleannode{thm:toric-distance}
\inputleannode{def:toric-code-with-distance}

\chapter{Bivariate-bicycle codes and the gross code}
\label{chap:bb}

Bivariate-bicycle codes are qLDPC codes built from two commuting polynomial
shift operators on an abelian group. The gross code is the $[[144,12,12]]$
member of the family; the argument here derives its distance from that of its
$[[72,12,6]]$ base by exhibiting the gross complex as a free $\Z_2$ cover of
the base and showing the distance doubles.

\inputleannode{def:x-double-cover}

\section{The base floor}

The base contributes a floor of six, established by ruling out every light
cycle with a finite certificate that the kernel checks directly.

\inputleannode{def:small-cycle-data}
\inputleannode{thm:small-cycle-floor}
\inputleannode{thm:base-distance}
\inputleannode{thm:gross-logical-weight}

\section{Doubling}

Pushing a gross logical down to the base gives weight at least six. Getting to
twelve requires knowing when the two sheets of the cover contribute
independently. They do outside a \emph{dangerous} sector, where the deck
transformation acts trivially on homology; there the light stabilizers are
classified explicitly and the surviving classes are confined to Smith cosets
whose floor is again twelve.

\inputleannode{thm:gross-sector-dichotomy}
\inputleannode{def:gross-stabilizer-code}
\inputleannode{thm:gross-distance}
\inputleannode{def:gross-code-with-distance}

\chapter{Code instances}
\label{chap:codes}

The remaining codes exercise the same machinery on smaller or differently
shaped examples. The $[[5,1,3]]$ code is the outlier: it is not CSS, and so the homological
route is unavailable, Its distance is settled instead by a direct search
over weight-one and weight-two anti-witnesses.

\inputleannode{def:steane7}
\inputleannode{def:shor9}
\inputleannode{def:five-qubit}
\inputleannode{def:four-qubit}
\inputleannode{def:iceberg}
\inputleannode{def:rotated-surface}
\inputleannode{def:repetition-n}
